\section*{الفصل \ref{ch:g12:deriv} --- الاشتقاق والتحدب}

\begin{solution}{exo:g12:deriv:1}
$f'(x) = 3x^2 - 5$.

بقاعدة خارج القسمة:
$g'(x) = \dfrac{(x^2+1) - x\cdot 2x}{(x^2+1)^2} = \dfrac{1 - x^2}{(x^2+1)^2}$.

بقاعدة السلسلة مع $u = 2x+1$: $h'(x) = 5 \cdot 2 \cdot (2x+1)^4 = 10(2x+1)^4$.

بقاعدة السلسلة مع $u = x^2 + 1$:
$k'(x) = \dfrac{2x}{2\sqrt{x^2+1}} = \dfrac{x}{\sqrt{x^2+1}}$.
\end{solution}

\begin{solution}{exo:g12:deriv:2}
$f'(x) = 2x - 3$، إذن $f(2) = -1$ و $f'(2) = 1$: \omterm{def:g12:deriv:tangent}{فالمماس} هو
$y = -1 + (x - 2) = x - 3$. والفجوة هي
\[
f(x) - (x - 3) = x^2 - 4x + 4 = (x-2)^2 \geq 0 ,
\]
إذن يقع \omterm{def:g10:functions:graph}{المنحنى} فوق \omterm{def:g12:deriv:tangent}{المماس} في كل مكان، ولا يلامسه إلا عند $x = 2$
(كما هو متوقع: \omterm{def:g11:func:function}{فالدالة} $f$ \omterm{def:g12:deriv:convex}{محدبة}).
\end{solution}

\begin{solution}{exo:g12:deriv:3}
كلاهما معدل تغير. فمع $f(x) = (1+x)^{100}$ عند $0$:
$f'(x) = 100(1+x)^{99}$، إذن النهاية هي $f'(0) = 100$.
ومع $g(x) = \sqrt{x}$ عند $4$: $g'(x) = \frac{1}{2\sqrt{x}}$، إذن النهاية هي
$g'(4) = \frac14$.
\end{solution}

\begin{solution}{exo:g12:deriv:4}
\omterm{def:g11:func:function}{مجموعة التعريف} $\R \setminus \{1\}$. والنهايات: عند $\pm\infty$، $f(x) \sim x \to
\pm\infty$؛ وعند $1^\pm$، البسط $\to 4 > 0$ والمقام $\to 0^\pm$،
إذن $f \to \pm\infty$؛ فالمستقيم $x = 1$ \omterm{def:g12:limcont:asymptote}{مقارب} شاقولي.

\[
f'(x) = \frac{2x(x-1) - (x^2+3)}{(x-1)^2} = \frac{x^2 - 2x - 3}{(x-1)^2}
= \frac{(x-3)(x+1)}{(x-1)^2}.
\]
ومنه $f' > 0$ على $\intoo{-\infty}{-1}$ وعلى $\intoo{3}{+\infty}$، و $f' < 0$ على
$\intoo{-1}{1}$ وعلى $\intoo{1}{3}$. فتتزايد \omterm{def:g11:func:function}{الدالة} إلى \omterm{def:g12:deriv:extremum}{قيمة عظمى محلية}
$f(-1) = -2$، ثم تتناقص إلى $-\infty$، ثم تقفز إلى $+\infty$ بعد \omterm{def:g12:limcont:asymptote}{المقارب}،
ثم تتناقص إلى \omterm{def:g10:functions:extrema}{قيمة صغرى} محلية $f(3) = 6$، ثم تتزايد.
\end{solution}

\begin{solution}{exo:g12:deriv:5}
من أجل $x \in \intoo{0}{15}$، تكون قاعدة العلبة مربعة طول ضلعها $30 - 2x$ و
ارتفاعها $x$، إذن
\[
V(x) = x(30 - 2x)^2 .
\]
$V'(x) = (30-2x)^2 + x \cdot 2(30-2x)(-2) = (30-2x)(30 - 2x - 4x)
= (30-2x)(30-6x)$. وعلى $\intoo{0}{15}$ يكون $30 - 2x > 0$، إذن تكون إشارة $V'$ هي إشارة
$30 - 6x$: موجبة قبل $x = 5$ وسالبة بعدها. فالحجم أعظمي من أجل
$x = 5$ cm، وقيمته $V(5) = 5 \times 20^2 = 2000\ \text{cm}^3$.
\end{solution}

\begin{solution}{exo:g12:deriv:6}
\emph{1.} $f'(x) = 4x^3 - 12x^2$ و $f''(x) = 12x^2 - 24x = 12x(x-2)$.
ومنه $f'' > 0$ على $\intoo{-\infty}{0}$ وعلى $\intoo{2}{+\infty}$ (\omterm{def:g12:deriv:convex}{محدبة})،
و $f'' < 0$ على $\intoo{0}{2}$ (\omterm{def:g12:deriv:convex}{مقعرة}). وتغيّر $f''$ إشارتها عند $0$ وعند $2$:
فكلتاهما \omterm{def:g12:deriv:inflection}{نقطة انعطاف}.

\emph{2.} عند $a = 2$: $f(2) = 16 - 32 + 10 = -6$، و $f'(2) = 32 - 48 = -16$،
\omterm{def:g12:deriv:tangent}{فالمماس} $y = -6 - 16(x-2) = -16x + 26$. والفجوة هي
\[
g(x) = f(x) + 16x - 26 = x^4 - 4x^3 + 16x - 16 .
\]
وبما أن $g(2) = g'(2) = g''(2) = 0$، فإن $(x-2)^3$ يقسم $g$؛ والقسمة تعطي
$g(x) = (x-2)^3(x+2)$. وبجوار $x = 2$ يكون $x + 2 > 0$، إذن تكون إشارة $g$ هي إشارة
$(x-2)^3$: سالبة قبل $2$ وموجبة بعدها. \omterm{def:g10:functions:graph}{فالمنحنى} يقطع
\omterm{def:g12:deriv:tangent}{المماس}: إذن $2$ \omterm{def:g12:deriv:inflection}{نقطة انعطاف} فعلًا.
\end{solution}

\begin{solution}{exo:g12:deriv:7}
\omterm{def:g11:func:function}{الدالة} $f(x) = \sqrt{x}$ تحقق
$f''(x) = -\frac{1}{4}x^{-3/2} < 0$ على $\intoo{0}{+\infty}$: فهي \omterm{def:g12:deriv:convex}{مقعرة}،
ومنه يقع منحناها \emph{تحت} كل \omterm{def:g12:deriv:tangent}{مماس} من مماساته. \omterm{def:g12:deriv:tangent}{والمماس} عند $a = 1$ هو
\[
y = f(1) + f'(1)(x - 1) = 1 + \tfrac12 (x-1) = \frac{x+1}{2},
\]
ومنه $\sqrt{x} \leq \frac{x+1}{2}$ من أجل كل $x > 0$. والمساواة تعني أن
\omterm{def:g10:functions:graph}{المنحنى} يلامس \omterm{def:g12:deriv:tangent}{المماس}، وهذا لا يحدث إلا عند نقطة التماس
$x = 1$. (وبصيغة مكافئة: $\left(\sqrt x - 1\right)^2 \geq 0$.)
\end{solution}

\begin{solution}{exo:g12:deriv:8}
حسب \cref{thm:g12:deriv:convexity}، يقع \omterm{def:g10:functions:graph}{منحنى} \omterm{def:g11:func:function}{دالة} \omterm{def:g12:deriv:convex}{محدبة} قابلة
للاشتقاق فوق كل \omterm{def:g12:deriv:tangent}{مماس} من مماساته. \omterm{def:g12:deriv:tangent}{والمماس} عند $a$ أفقي
($f'(a) = 0$)، ومعادلته $y = f(a)$. ومنه $f(x) \geq f(a)$ من أجل كل
$x \in \R$: أي إن $f(a)$ هي \omterm{def:g10:functions:extrema}{القيمة الصغرى} الشاملة.
\end{solution}

\begin{solution}{exo:g12:deriv:9}
\emph{المحيط الثابت.} مستطيل ضلعاه $x$ و $\frac{p}{2} - x$
($0 < x < \frac p2$) مساحته $S(x) = x\left(\frac p2 - x\right)$. ثم
$S'(x) = \frac p2 - 2x$ تنعدم عند $x = \frac p4$، وهي موجبة قبلها و
سالبة بعدها: \omterm{def:g10:functions:extrema}{فالقيمة العظمى} عند $x = \frac{p}{4}$، حيث يتساوى الضلعان
عند $\frac p4$ --- أي مربع.

\emph{المساحة الثابتة.} الضلعان $x$ و $\frac Ax$ يعطيان محيطًا
$P(x) = 2\left(x + \frac Ax\right)$، مع
$P'(x) = 2\left(1 - \frac{A}{x^2}\right)$، وهي سالبة من أجل $x < \sqrt A$ و
موجبة بعدها: \omterm{def:g10:functions:extrema}{فالقيمة الصغرى} عند $x = \sqrt{A}$، أي مربع من جديد.

\emph{العلاقة.} العبارتان متقابلتان. لنفترض أن مستطيلًا غير مربع
يصغّر المحيط عند مساحة ثابتة $A$؛ فالمربع ذو المحيط نفسه
ستكون مساحته $> A$ حسب العبارة الأولى، وتصغيره تحاكيًا
حتى المساحة $A$ سينقص محيطه تمامًا --- وهذا يناقض
كونه أصغريًا. فكل عبارة تستلزم الأخرى إذن.
\end{solution}

\begin{solution}{pb:g12:deriv:1}
\textbf{1.} $30x\left(3x^2 + 1\right)^4$؛ \quad
$\dfrac{x}{\sqrt{x^2 + 9}}$؛ \quad
$-\dfrac{2x}{\left(x^2 + 1\right)^2}$.

\textbf{2.} $y(4) = 5$، \omterm{def:g10:reffunc:affine}{والميل} $\frac45$: \omterm{def:g12:deriv:tangent}{فالمماس}
$y = 5 + \frac45(x - 4)$، أي $y = \frac45 x + \frac95$.

\textbf{3.} $f'(x) = \frac{(x^2 + 1) - x \cdot 2x}{(x^2+1)^2}
= \frac{1 - x^2}{(x^2 + 1)^2}$: سالبة ثم موجبة ثم سالبة
عبر $-1$ و $1$: \omterm{def:g10:functions:extrema}{فالقيمة الصغرى} $f(-1) = -\frac12$، \omterm{def:g10:functions:extrema}{والقيمة العظمى}
$f(1) = \frac12$، مع \omterm{def:g12:limcont:asymptote}{المقارب} الأفقي $0$ على الجهتين.

\textbf{4.} $g'(x) = 3x^2 - 6x = 3x(x - 2)$: \omterm{def:g12:seq:monotonic}{متزايدة} ثم
\omterm{def:g10:functions:variations}{متناقصة} ثم \omterm{def:g12:seq:monotonic}{متزايدة} عبر $0$ و $2$؛ \omterm{def:g12:deriv:extremum}{فالقيمة العظمى المحلية}
$g(0) = 4$، والصغرى المحلية $g(2) = 0$. و $g''(x) = 6x - 6$: \omterm{def:g12:deriv:convex}{مقعرة}
قبل $1$ \omterm{def:g12:deriv:convex}{ومحدبة} بعدها: فالانعطاف عند $(1, 2)$.

\textbf{5.} \omterm{def:g12:deriv:tangent}{المماس} عند $1$: \omterm{def:g10:reffunc:affine}{الميل} $g'(1) = -3$:
$y = 2 - 3(x - 1) = -3x + 5$. والفرق:
$x^3 - 3x^2 + 4 - (-3x + 5) = x^3 - 3x^2 + 3x - 1 =
(x - 1)^3$، وهو يغيّر إشارته عند $1$: \omterm{def:g12:deriv:tangent}{فالمماس}
\emph{يقطع} \omterm{def:g10:functions:graph}{المنحنى} --- وهو السلوك \omterm{def:g11:quad:discriminant}{المميز} عند
\omterm{def:g12:deriv:inflection}{نقطة انعطاف}، حيث يبدّل \omterm{def:g10:functions:graph}{المنحنى} جهته.

\textbf{6.} $T(x) = \dfrac{\sqrt{x^2 + 900}}{5} +
\dfrac{\sqrt{(40 - x)^2 + 400}}{2}$. فالدخول قبل $0$ أو
بعد $40$ يطيل \emph{المسارين} معًا: أي بلا جدوى.

\textbf{7.} $T'(x) = \dfrac{x}{5\sqrt{x^2 + 900}} -
\dfrac{40 - x}{2\sqrt{(40 - x)^2 + 400}}$.

\textbf{8.} في مثلث الركض، يكون الضلع على طول الشاطئ
هو $x$ والوتر $\sqrt{x^2 + 900}$: \omterm{def:g11:trigo:cossin}{فجيب}
الزاوية مع العمود هو بالضبط خارج قسمتهما
(المقابل على الوتر)؛ وبالمثل
$\sin\theta_2 = \frac{40 - x}{\sqrt{(40-x)^2 + 400}}$. عندئذٍ
يُقرأ $T'(x) = 0$ على الصورة
$\frac{\sin\theta_1}{5} = \frac{\sin\theta_2}{2}$ --- وهو قانون
سنيل بالسرعتين $5$ و $2$.

\textbf{9.} $T(24) \approx 20.5$ s؛ و $T(40) = 10 + 10 = 20$ s؛
و $T(0) = 6 + \sqrt{2000}/2 \approx 28.4$ s. فالخط المستقيم
يخسر حتى أمام ``اركض إلى النهاية ثم اسبح مباشرة'' ---
وكلاهما يخسر أمام المسار المنكسر.

\textbf{10.} التنصيف على $T'$ (وهي سالبة عند $24$ وموجبة عند
$40$) يتقارب إلى $x^* \approx 34$ m، مع
$T(x^*) \approx 19.5$ s: أي أسرع بنحو ثانية واحدة من
الخط المستقيم --- وهو الفرق بين إنقاذ و
مأساة.

\textbf{11.} زمن كل مسار \omterm{def:g11:func:function}{دالة} \omterm{def:g12:deriv:convex}{محدبة} في $x$
(وهو فرع $\sqrt{\text{التربيعي}}$)، إذن $T$ \omterm{def:g12:deriv:convex}{محدبة}، والنقطة
الحرجة \omterm{def:g11:func:function}{لدالة} \omterm{def:g12:deriv:convex}{محدبة} قابلة للاشتقاق هي قيمتها
الصغرى الشاملة (\cref{exo:g12:deriv:8}): أي إن $x^*$ ليست
مستقرة فحسب، بل لا تُقهر.

\textbf{12.} الضوء في الماء أبطأ؛ وحسب مبدأ فيرما
ينكسر أسرع مسار من الهواء إلى الماء عند السطح بالضبط
حسب قانون سنيل --- فتصل الأشعة من قلم مغمور إلى العين
منكسرة، ويرى الدماغ، وهو يمدّ الخطوط المستقيمة، القلم
مكسورًا عند خط الماء.

\textbf{13.} $(x^2)'' = 2 > 0$: إذن \omterm{def:g12:deriv:convex}{محدبة}. \omterm{def:g12:deriv:convex}{والتحدب} عند
\omterm{prop:g10:coordgeom:midpoint}{المنتصف}: $f\left(\frac{a+b}{2}\right) \leq
\frac{f(a) + f(b)}{2}$، وهذا هو المطلوب. وباليد:
$\frac{a^2 + b^2}{2} - \left(\frac{a+b}{2}\right)^2 =
\frac{(a - b)^2}{4} \geq 0$.

\textbf{14.} بأخذ الجذور التربيعية (وهي \omterm{def:g12:seq:monotonic}{متزايدة}) في السؤال 13:
$\frac{a + b}{2} \leq \sqrt{\frac{a^2 + b^2}{2}}$. والسلسلة الكاملة
على $2, 8$: التوافقي $\frac{2 \cdot 16}{10} = 3.2$؛ والهندسي
$\sqrt{16} = 4$؛ والحسابي $5$؛ والتربيعي
$\sqrt{34} \approx 5.83$: فكل \omterm{def:g10:stats:mean}{متوسط} ينحني للذي يليه.

\textbf{15.} \omterm{def:g12:deriv:tangent}{المماس} \omterm{def:g10:functions:graph}{للمنحنى} $\frac1x$ عند $1$: القيمة $1$ \omterm{def:g10:reffunc:affine}{والميل}
$-1$: أي $y = 2 - x$. والمنحنيات \omterm{def:g12:deriv:convex}{المحدبة} تجلس فوق مماساتها
(\cref{thm:g12:deriv:convexity}): إذن $\frac1x \geq 2 - x$ من أجل كل
$x > 0$، مع المساواة عند نقطة التماس $x = 1$ بالضبط.

\textbf{16.} عند الانعطاف، يبلغ العدد \emph{اليومي} للحالات
(وهو \omterm{def:g12:deriv:derivative}{المشتقة}) ذروته: فيتوقف النمو عن التسارع ويبدأ
بالتباطؤ --- وهي أول إشارة رياضية على أن الموجة
تنقلب، قبل أن تنخفض الأعداد نفسها بزمن طويل. فقبلها،
$f'' > 0$ (كل يوم أسوأ من سابقه)؛ وبعدها،
$f'' < 0$ (ما زال النمو قائمًا، لكن أبطأ).

\textbf{17.} $S'(r) = 4\pi r - \frac{2V}{r^2} = 0$ يعطي
$r^3 = \frac{V}{2\pi} = \frac{330}{2\pi}$:
أي $r \approx 3.74$ cm، و
$h = \frac{330}{\pi r^2} \approx 7.49$ cm.

\textbf{18.} $h = \frac{V}{\pi r^2} = \frac{2\pi r^3}{\pi r^2}
= 2r$ عند الأمثلية: فالارتفاع يساوي القطر، مهما كان
الحجم.

\textbf{19.} $S''(r) = 4\pi + \frac{4V}{r^3} > 0$: إذن \omterm{def:g12:deriv:convex}{محدبة}، ومنه يكون
نصف القطر الحرج هو \omterm{def:g10:functions:extrema}{القيمة الصغرى} الشاملة. أما العلب الحقيقية فتقف
أطول من عرضها لأجل المسك والتكديس وحضور الرف والمعدن الأثخن
في الأعلى والأسفل --- فالأمثلة تصغّر دائمًا بالضبط
ما كتبتَه، لا ما قصدتَه.

\textbf{20.} المنقذة: انمذج $T(x)$، ثم اشتق، ثم سنيل،
ثم \omterm{def:g12:deriv:convex}{التحدب}، فثانية واحدة موفَّرة. والعلبة: انمذج $S(r)$، ثم اشتق،
فيكون $h = 2r$، ثم \omterm{def:g12:deriv:convex}{التحدب}، فمعدن موفَّر. والعارضة
(\cref{pb:g11:deriv:1}): انمذج $S(w)$، ثم اشتق،
فيكون $d = w\sqrt2$، ثم الجدول، فصلابة مكتسبة. وحسب مبدأ
فيرما، يشغّل الضوء نفسه هذا الخط عند كل سطح
يلقاه --- فالطبيعة كانت أول من أَمثَل.
\end{solution}
